\section{¿De dónde vienen las recompensas?}

En las clases anteriores, supusimos que la función de recompensa $r(s, a)$ estaba dada. Pero en la realidad, \textbf{definir la función de recompensa es en sí un problema difícil}.

\begin{knowledgebox}{Situación de la recompensa en distintos dominios}
\begin{itemize}
    \item \textbf{Videojuegos}: la puntuación es naturalmente la recompensa, lista para usarse directamente.
    \item \textbf{Robótica}: ¿qué es un «buen agarre»? ¿qué es un «doblado ordenado»? Es necesario definirlo manualmente, y a menudo se usan métricas sustitutas.
    \item \textbf{Sistemas de diálogo}: ¿qué es una «buena respuesta»? Es altamente subjetivo y difícil de medir con una función simple.
    \item \textbf{Conducción autónoma}: implica compensaciones entre múltiples dimensiones como la seguridad, la comodidad y la eficiencia.
\end{itemize}
\end{knowledgebox}

Además de definir la recompensa directamente, ya vimos una alternativa: \textbf{el aprendizaje por imitación}, que imita directamente las acciones del experto sin necesidad de una recompensa. Pero el aprendizaje por imitación no razona sobre los resultados ni la dinámica, y el experto puede tener un embodiment distinto.

\subsection{Resumen de la sección}

El diseño de recompensas es un cuello de botella clave en las aplicaciones de RL. Esta clase presenta tres métodos para \textbf{aprender} la recompensa a partir de la supervisión humana.

